% !TeX root = ../../tesis.tex
%%%%%%%%%%%%%%%%%%%%%%%%%%%%%%%%%%%%%%%%%%%%%%%%%%%%%%%%%%%%%%%%%%%
% §2.1.1 VULNERABILIDAD Y ROBUSTEZ DE REDES DE TRANSPORTE — EN BLANCO
% Origen: 3-Conceptos-Indicadores/3-1-Definiciones/3-1-2-Vulnerabilidad.tex
%   (98 l.; §3.1.2 completo, pedido explícito del asesor):
%   - l. 8-31  «Concepto y Marco Teórico desde las Redes Complejas»
%   - l. 32-98 «Fragilidad Estructural y Análisis Topológico», con la
%     figura fig:topologias-red (Lotero et al.) y su cita en el texto.
% Autores: lotero2014vulnerabilidad, albert2000error,
%   holmgren2006vulnerability, nagurney2009fragile, ouyang2009vulnerability.
% Al mover: quitar los títulos \subsubsection (texto continuo, acuerdo 4)
%   y conservar \label{subsec:...} que se citan desde otros capítulos.
% En §3.1 queda solo la definición operativa breve que necesite el Cap. 3.
%%%%%%%%%%%%%%%%%%%%%%%%%%%%%%%%%%%%%%%%%%%%%%%%%%%%%%%%%%%%%%%%%%%
